\def\topic{many-valued}
\ifdefined\MyHeaderIncluded
\else
\def\qwertyflag{true}

\ifdefined\MyHeaderIncluded

\else
  \newcommand{\foobar}{FooBar}
  \def\MyHeaderIncluded{true}
% if is till the end of the file

\documentclass[10pt,a4paper]{article}
\usepackage[utf8]{inputenc}
%\usepackage[german]{babel}
%\usepackage[cyrillic]{babel}
\usepackage[T1]{fontenc}
%\usepackage[T2A]{fontenc}
%\usepackage[english, russian]{babel}
\usepackage{amsmath}
\usepackage{amsfonts}
\usepackage{amssymb}
\usepackage{amsthm}
\usepackage{mathtools} %coloneqq

\usepackage{bussproofs} %

%\renewcommand{\thesubsection}{\arabic{subsection}}
%\newtheorem{thm}{Theorem}[section]
%\renewcommand{\thethm}{\arabic{subsection}.\arabic{thm}}

\usepackage{thmtools}
\usepackage{hyperref}
\hypersetup{
    colorlinks=true,
    linkcolor=blue,
    filecolor=magenta,      
    urlcolor=cyan,
    pdftitle={Overleaf Example},
    %pdfpagemode=FullScreen,
    }
%\usepackage{stix-type1}%\turnediota
\theoremstyle{definition}

\newtheorem{constructor}{Constructor}[section]
\newtheorem{iaxiom}{Axiom}
\newtheorem{axiom}{Axiom}
\newtheorem*{axiom*}{Axiom}
\newtheorem{definition}{Definition}[section]
\newtheorem*{definition*}{Definition}
\newtheorem{remark}{Remark}[section]
\newtheorem*{remark*}{Remark}
\newtheorem{recall}{Recall}[section]
\newtheorem*{recall*}{Recall}

\newtheorem{lemma}{Lemma}[section]
\newtheorem{theorem}{Theorem}[section]
\newtheorem*{theorem*}{Theorem}
\newtheorem{note}{Note}[section]
\usepackage{etoolbox}

%\newcommand*{\@Dashv}[2]{%
%  \reflectbox{$\m@th#1#2$}%
%}
\newcommand*{\separation}[3]{% 1)binder, 2)set, 3)condition
  \{#1\in #2: #3\}
}

% https://www.freecodecamp.org/news/simple-regex-tricks-for-beginners-3acb3fa257cb/
%  \\mathcal\{T\} |---> \topology
%  \\mathcal\{([^\}]*)\}   |--->  $1
\newcommand{\AND}{\ \ \&\ \ }
\newcommand{\any}{\mathrm{\ any}}
\newcommand{\borel}{\mathcal{B}}
\newcommand{\can}{\mathrm{can}}
\newcommand{\canonize}{\mathrm{canonize}}
\newcommand{\closure}{\mbox{cl}}
\newcommand{\complex}{\mathbb{C}}
\newcommand{\converges}{\longrightarrow}
\newcommand{\countable}{\mbox{\ countable}}
\newcommand{\ldefined}{\mbox{\ defined}}
\newcommand{\dom}{\mbox{dom}}
\newcommand{\Dom}{\mbox{Dom}}
\newcommand{\Done}{\noindent \textbf{Done.}\\}
\newcommand{\lexists}{\mbox{\ exists}}
\newcommand{\lforall}{\mbox{\ forall}}
\newcommand{\lunique}{\mbox{\ unique}}
\newcommand{\family}{\mathrm{\ family}}
\newcommand{\Fm}{\mathrm{Fm}}
\newcommand{\FuncSymb}{\mathrm{FuncSymb}}
\newcommand{\function}{\mathrm{\ function}}
\newcommand{\functional}{\mbox{\ functional}}
\newcommand{\IFF}{\ \Longleftrightarrow\ }
\renewcommand{\iff}{\ \longleftrightarrow\ }
\newcommand{\IMPLIES}{\ \Longrightarrow\ }
\renewcommand{\implies}{\ \longrightarrow\ }
\newcommand{\Inf}{\mathrm{Inf}}
\newcommand{\injective}{\mbox{\ injective}}
\newcommand{\injection}{\mbox{\ injection}}
\newcommand{\interior}{\mbox{int}}
\newcommand{\intersection}{\bigcap}
\newcommand{\isempty}{\mathrm{\ empty}}
\newcommand{\LB}{\mathrm{LB}}
\newcommand{\nghbset}{\mathcal{N}}
\newcommand{\nonempty}{\mathrm{\ nonempty}}
\newcommand{\Ord}{\mathrm{Ord}}
\newcommand{\power}{\mathcal{P}}
\newcommand{\pr}{\mathrm{pr}}
\newcommand{\fst}{\mathrm{fst}}
\newcommand{\snd}{\mathrm{snd}}
\newcommand{\PredSymb}{\mathrm{PredSymb}}
\newcommand{\Proof}{\noindent \textbf{Proof.}\\}
\newcommand{\ran}{\mathrm{ran}}
\newcommand{\Ran}{\mathrm{Ran}}
\newcommand{\real}{\mathbb{R}}
\newcommand{\relation}{\mathrm{\ relation}}
\newcommand{\set}{\mathrm{\ set}}
\newcommand{\singleton}{\mathrm{\ singleton}}
\newcommand{\semantics}{\mathrm{semantics}}
\newcommand{\semialgebra}{\mathcal{S}}
\newcommand{\Sup}{\mathrm{Sup}}
\newcommand{\surjective}{\mbox{\ surjective\ on\ }}
\newcommand{\Term}{\mathrm{Term}}
\newcommand{\Type}{\mathrm{Type}}
\newcommand{\topology}{\mathcal{T}}
\newcommand{\total}{\mbox{\ total\ on\ }}
\newcommand{\transitive}{\mathrm{\ transitive}}
\newcommand{\union}{\bigcup}
\newcommand{\lTrue}{\mathrm{True}}
\newcommand{\lFalse}{\mathrm{False}}
\newcommand{\lUncertain}{\mathrm{Uncertain}}
\newcommand{\lif}{\mathrm{if}}
\newcommand{\lthen}{\mathrm{then}}
\newcommand{\lelse}{\mathrm{else}}
\newcommand{\lexception}{\mathrm{exception}}
\newcommand{\lundef}{\mathrm{undef}}
\newcommand{\univocal}{\mathrm{\ univocal}}
\newcommand{\univocalat}{\mathrm{\ univocal\ at\ }}
\newcommand{\UB}{\mathrm{UB}}
\newcommand{\Un}{\mathcal{U}}
\newcommand{\Var}{\mathrm{Var}}

\newcommand{\notvarphi}{\phi}
\renewcommand{\phi}{\varphi}

\newcommand{\comment}[1]{\quad\mbox{(#1)}}
% local commands
%\preto\section{\undef\blipo} %etoolbox
\preto\subsection{\undef\blipo} %etoolbox
% https://tex.stackexchange.com/questions/105684/limit-macro-newcommand-scope-to-section
\usepackage{graphicx}
\usepackage{authblk}

\setlength\parindent{0pt} % whole doc of \noindent
\usepackage{soul}
%\usepackage{mathgl}

\fi % guard
\ifdefined\topic
\else
\def\topic{No Titel}
\fi
\title{\topic}
\author[1]{Georgy Dunaev}
\begin{document}
\maketitle
\fi
\section{Many-valued ideas}

To denote the constructor of an uncertain formula, we will use the symbol $\uncert$, because it is symmetrical (neither $\bot$, nor $\top$ are  symmetrical ). And it looks like crossed D, and "defined" starts with D, and it is connected with not being defined: normally(but not always!), if $f$ is a functional symbol, $f(...,\lundef,...)$.\\
\[
I_\Term(\uncert) \coloneqq U
\]

============

Remark: actually we can use 6 truth values.\\
 True = (present, stable(True))\\
False = (present, stable(False))\\
U1    = ( absent, stable(True))\\
U2    = ( absent, stable(False))\\
U3    = (present, unstable)\\
U4    = ( absent, unstable)\\
Connectives are contagious in the first argument, it means the following:\\
Let's say \# is an interpretation of a connective, and v is a truth value, then\\
 (absent, ...) \# v = (absent, ...)\\
 (note, that "..." in lhs and rhs are not necessary the same)\\
 (..., stable(True)) $\land$ (..., stable(True)) =  (..., stable(True))\\
But two predicates can stabilize each other, so these truth values does not have a truth table, that's why we will not use them.\\

Also therefore Uncertain OR Uncertain is not necessary Uncertain in such semantics!
Therefore we need to require even more: the certainity of every subformula, otherwise it will complicate the system too much?\\

e.g. $(f(x)=5) \lor (x\notin\dom(f))$.\\
Or will it require too much to require?\\
Yes, it will be too much, since then\\
$(f\function)\land(x\in\dom(f)) \implies (f(x)=f(x))$\\
will not work normally.\\
Therefore we need something more interesting like:\\
 $(\lFalse\implies\lUncertain) := \lTrue$.\\
Or even brave \& better: reject the whole idea of having a defined truth table for uncertain values(or at least having a full/complete truth table).
Such idea was wrong, because it is not true even in completely classical logic of existence: the formula "$(x\in y)$" is not valid, so it has no any particular defined truth value.\\

=================


============================================\\
Idea: Truth values of predicates are subconglomerates of $(2\times)\Un\times\dots\times\Un$ of values, where $I(P)$ is true.\\
"(Always)True" corresponds to the whole subconglomerate $(1\times)\Un\times\dots\times\Un$.\\
"(Always)False" corresponds to the empty subconglomerate $(0\times)\emptyset\times\dots\times\emptyset$, which is simply $\emptyset$.\\
"Uncertain" corresponds to any other subconglomerate.\\
All products here are given $n$ times: $n$ is a valency of the predicate symbol $P$.
============================================\\

For example, the following predicates:\\
\begin{itemize}
\item (of valency 0):\\
"$\bot$" corresponds to 0.("False")\\
"$\top$" corresponds to 1.("True")\\
"$\uncert$" corresponds to 2.("Uncertain")\\
(from now on, we will forget about "$2\times$" form the definition, for simplicity)
\item (of valency 1):\\
"(X unique)" corresponds to the collection of all conglomerate-semisingletones.\\
"(X exists)" corresponds to the collection of all nonempty conglomerates.\\
"(X (well-)defined)" corresponds to the collection of all nonempty conglomerate-semisingletones, i.e. to the collection of all conglomerate-singletones\\
"empty" \ corresponds to the collection of empty conglomerates, i.e. just $\{\emptyset\}$.
\item (of valency 2):\\
"$(x\in y)$" corresponds to collection of all ordered pairs of nonempty conglomerates, where each class from first conglomerate is in each class of the second conglomerate.\\
"$(x\notin y)$" corresponds to collection of all ordered pairs of nonempty conglomerates, where each class from first conglomerate is not in any class of the second conglomerate.
\end{itemize}
\begin{definition}(Theorem)
Theorem is a predicate which has the truth value "True".\\
\end{definition}
%$\tau t$
%Therefore.
%This is something unresolved, we need some time...


% and define the Uncertain value, like:\\
%Uncertain = Non-existent OR Non-constant\\
%\,\quad\quad = Absent OR Unstable\\
%\,\quad\quad = Void OR Varying\\
%But then, how connectives will work? They will be "contagious" take "the worst case"? %Not necessary: two predicates can stabilize each other.\\
Idea 2:\\

Actually, variables are ordered: $x_1,x_2,x_3,\dots$.
So we can consider the following correspondence evaluation of the formula $x_2\in x_4$ via assignment function $w$, with 
\[
(1\times)\Un\times \{w(x_2)\}\times\Un\times\{w(x_4)\}\times\Un\times\Un\times\Un\dots,
\]
where "infinite product" is abuse of notation, which means a collection of functions $f$ from $\mathbb{N}_0$ to $\power\Un$, such that $f(0)=1$. %\subseteq 

\iffalse
\symbol{200}\symbol{201}\symbol{202}\symbol{203}
\symbol{204}\symbol{205}\symbol{206}\symbol{207}
\symbol{208}\symbol{209}\symbol{210}\symbol{211}
\symbol{212}\symbol{213}\symbol{214}\symbol{215}
\symbol{216}\symbol{217}\symbol{218}\symbol{219}\\
\symbol{220}\symbol{221}\symbol{222}\symbol{223}
\symbol{224}\symbol{225}\symbol{226}\symbol{227}
\symbol{228}\symbol{229}\symbol{230}\symbol{231}
\symbol{232}\symbol{233}\symbol{234}\symbol{235}
\symbol{236}\symbol{237}\symbol{238}\symbol{239}\\
\symbol{240}\symbol{241}\symbol{242}\symbol{243}
\symbol{244}\symbol{245}\symbol{246}\symbol{247}
\symbol{248}\symbol{249}\symbol{250}\symbol{251}
\symbol{252}\symbol{253}\symbol{254}\symbol{255}
\fi

%\symbol{253A}


\ifdefined\qwertyflag
\undef\qwertyflag
\end{document}
\fi